\chapter{Análisis de impacto}

\section{Impacto Personal}
El desarrollo de este TFG ha supuesto un impacto personal significativo en términos de adquisición de conocimientos profundos sobre la teoría de polinomios ortogonales de Sobolev, el manejo avanzado de herramientas de cálculo simbólico como \textsc{Maple}, y el desarrollo de habilidades de investigación, análisis crítico y comunicación científica.

\section{Impacto en Investigación}
En el ámbito de la investigación, este trabajo busca:
\begin{itemize}
	\item \textbf{Aportar nuevas herramientas computacionales y de visualización} que permitan la experimentación numérica y la formulación de nuevas conjeturas en el campo.
	\item \textbf{Ofrecer evidencia empírica y vías futuras} sobre el comportamiento de los ceros y los valores singulares en casos poco estudiados, como aquellos donde no hay dominación secuencial o en presencia de átomos discretos.
\end{itemize}

\section{Beneficios Esperados y Limitaciones}
Los resultados y herramientas desarrolladas en este TFG ofrecen una valiosa contribución a la teoría matemática fundamental, conectando con problemas abiertos en teoría espectral de matrices infinitas y teoría del potencial en el plano complejo. Además, refuerzan la comprensión de fenómenos de transición de fase en matrices aleatorias con restricciones derivativas, con aplicaciones relevantes en física estadística y teoría cuántica de la información, como se menciona en la introducción.

No obstante, es importante reconocer ciertas limitaciones, principalmente relacionadas con la complejidad computacional. El análisis de polinomios de grado alto y la manipulación de matrices de momentos pueden demandar recursos computacionales significativos, lo que podría restringir algunas aplicaciones prácticas o requerir optimizaciones adicionales para su implementación efectiva.

\section{Impacto en los Objetivos de Desarrollo Sostenible (ODS)}
Este trabajo, aunque fundamentalmente teórico y computacional en el ámbito de las matemáticas, puede contribuir indirectamente a ciertos Objetivos de Desarrollo Sostenible:

\begin{itemize}
	\item \textbf{ODS 4: Educación de Calidad.}
	      Las herramientas y visualizaciones desarrolladas pueden enriquecer la educación superior en matemáticas, sirviendo como material didáctico avanzado y apoyando la formación científica.

	\item \textbf{ODS 9: Industria, Innovación e Infraestructura.}
	      El estudio de polinomios ortogonales de Sobolev puede impulsar innovaciones en modelos numéricos para problemas industriales complejos, contribuyendo a la infraestructura científica necesaria para el avance tecnológico.

	\item \textbf{ODS 17: Alianzas para Lograr los Objetivos.}
	      La publicación de resultados y código abierto promueve la colaboración científica internacional, facilitando la construcción colectiva de conocimiento y el fortalecimiento de redes académicas.
\end{itemize}

